\documentclass[master, fontsize=12pt]{diploma}

\usepackage[supervisor]{review}

\student{Инютин Максим Андреевич}
\faculty{№ 8 <<Компьютерные науки и прикладная математика>>}
\department{Образовательный центр Института № 8}
\speciality{01.04.02 Прикладная математика и информатика}
\profile{Продуктовая разработка и разработка программных систем}
\group{М8О-206СВ-24}

\theme{Разработка алгоритма лидарной одометрии беспилотного автомобиля, учитывающего перекрытие поля зрения}
% Дата. Оставляем пустое место для дня
\workDate{\uline{\hspace{24pt}} мая \the\year\ года}

\supervisor{Ухов Пётр Александрович}
\supervisorCredentials{кандидат технических наук, доцент, преподаватель кафедры 806}

\firstConsultant{Евсеев Дмитрий Сергеевич}
\firstConsultantCredentials{доктор физико-математических наук, профессор, профессор кафедры 806}

% \secondConsultant{Алешина Елена Константиновна}
% \secondConsultantCredentials{доктор физико-математических наук, профессор, профессор кафедры 806}

% \reviewer{Свиридов Михаил Артемьевич}
% \reviewerCredentials{доктор физико-математических наук, профессор}

\headOfDepartmentTitle{И. о. директора Образовательного центра}
\headOfDepartment{Крылов Сергей Сергеевич}



\begin{document}
\fillReviewHeader

Выпускная квалификационная работа Инютина Максима Андреевича посвящена разработке алгоритма лидарной одометрии беспилотного автомобиля, учитывающего перекрытие поля зрения. Тема является актуальной для современных систем автономного вождения, поскольку качество оценки собственного движения непосредственно влияет на локализацию, планирование и безопасность движения, а крупные динамические объекты могут существенно ухудшать наблюдаемость сцены.

В ходе выполнения работы студент проявил самостоятельность, системность и хорошую инженерную культуру. Были изучены существующие подходы к регистрации облаков точек и лидарной одометрии, рассмотрены ограничения классического ICP в условиях динамического перекрытия поля зрения, описан набор данных Boreas-RT и подготовлен воспроизводимый сценарий вычислительного эксперимента. Автор разработал программный конвейер обработки последовательности лидарных кадров, включающий предобработку данных, двойную вокселизацию, взвешенный ICP, оценку преобразования между облаками точек, модель постоянной скорости и фильтр Калмана.

Отдельно следует отметить, что экспериментальная часть работы выполнена не как демонстрация на одном удачном примере, а как серия сравнений на 725 десятисекундных сценах из 10 тестовых проездов. Автор анализирует как успешные случаи, где фильтр Калмана стабилизирует начальное приближение и уменьшает деградацию оценки продольной скорости, так и ограничения метода в геометрически сложных сценах. Такой подход делает выводы работы достаточно взвешенными и практически полезными.

\makeSimilarityLine

Заключение: выпускная квалификационная работа выполнена в полном объёме, соответствует требованиям, предъявляемым к магистерским диссертациям, и заслуживает положительной оценки, а Инютин Максим Андреевич заслуживает присвоения квалификации магистра по направлению подготовки 01.04.02 <<Прикладная математика и информатика>>.

\makeReviewSignature
\end{document}